\documentclass[12pt,a4paper]{article}
\usepackage[utf8]{inputenc}
\usepackage[T1]{fontenc}
\usepackage{amsmath,amssymb,amsthm}
\usepackage{geometry}
\usepackage{enumitem}
\usepackage{hyperref}
\usepackage{fancyhdr}
\usepackage{titlesec}

\geometry{margin=1in}
\pagestyle{fancy}
\fancyhf{}
\rhead{Lecture 5}
\lhead{Quantum Systems via Linear Algebra}
\cfoot{\thepage}

\titleformat{\section}{\large\bfseries}{}{0em}{}
\titleformat{\subsection}{\normalsize\bfseries}{}{0em}{}

\newtheorem{theorem}{Theorem}
\newtheorem{definition}{Definition}

\title{\textbf{Lecture 5: Commutators and Two-Level Dynamics} \\ \large Simultaneous Measurability, Non-Commutativity, and Precession}
\author{Quantum Systems via Linear Algebra, Differential Equations, and Probability}
\date{}

\begin{document}
\maketitle
\thispagestyle{fancy}

% ============================================================
\section{1. Simultaneous Measurability and Commutators}
% ============================================================

\subsection{1.1 When can two observables be measured simultaneously?}

When you measure observable $L$, you leave the system in an eigenstate of $L$.  When you measure a second observable $M$, you leave it in an eigenstate of $M$.  If both measurements can be performed in either order without interference, the system must end up in a state that is an eigenvector of \emph{both} matrices.  This requires a complete basis of \textbf{simultaneous eigenvectors}.

\begin{theorem}
Two self-adjoint matrices $L$ and $M$ can be simultaneously measured if and only if they \textbf{commute}:
\[
    \boxed{[L,M] \equiv LM - ML = 0.}
\]
\end{theorem}

\subsection{1.2 Proof}
$(\Rightarrow)$ Suppose $\{\mathbf{e}_i\}$ is a complete set of simultaneous eigenvectors:
\[
    L\mathbf{e}_i = l_i\mathbf{e}_i, \qquad M\mathbf{e}_i = m_i\mathbf{e}_i.
\]
Apply $ML$ and $LM$ to any basis vector:
\[
    ML\mathbf{e}_i = M(l_i\mathbf{e}_i) = l_i m_i\mathbf{e}_i, \qquad
    LM\mathbf{e}_i = L(m_i\mathbf{e}_i) = m_i l_i\mathbf{e}_i.
\]
Since $l_i$ and $m_i$ are just numbers, $l_i m_i = m_i l_i$.  So $ML\mathbf{e}_i=LM\mathbf{e}_i$ for every basis vector.  Any vector is a superposition of basis vectors, so $ML=LM$ on the entire space.

\medskip
\noindent$(\Leftarrow)$ The converse also holds: if $[L,M]=0$ then a common eigenbasis exists.

\subsection{1.3 The uncertainty idea}
When $[L,M]\neq 0$, there are \emph{no} simultaneous eigenstates.  If you measure $L$ and leave the system in an eigenstate of $L$, that state is \emph{not} an eigenstate of $M$.  A subsequent measurement of $M$ will give a statistical result, \emph{and} it will disturb the value of $L$.  This is the origin of the \textbf{uncertainty principle}: certain pairs of observables cannot both have definite values at the same time.

\medskip
\noindent\textbf{Engineering analog:} Non-commuting observables are analogous to non-commuting transfer functions in feedback loops---the order of operations matters.

\subsection{1.4 Commutation relations for $X$, $Y$, $Z$}
The three standard observable matrices satisfy the cyclic commutation relations:
\[
    \boxed{[Z,X] = 2iY, \qquad
    [X,Y] = 2iZ, \qquad
    [Y,Z] = 2iX.}
\]
These can be verified by explicit matrix multiplication.  For example:
\[
    ZX = \begin{pmatrix}1&0\\0&-1\end{pmatrix}\begin{pmatrix}0&1\\1&0\end{pmatrix} = \begin{pmatrix}0&1\\-1&0\end{pmatrix},
\]
\[
    XZ = \begin{pmatrix}0&1\\1&0\end{pmatrix}\begin{pmatrix}1&0\\0&-1\end{pmatrix} = \begin{pmatrix}0&-1\\1&0\end{pmatrix},
\]
\[
    [Z,X] = \begin{pmatrix}0&2\\-2&0\end{pmatrix} = 2i\begin{pmatrix}0&-i\\i&0\end{pmatrix} = 2iY.\;\checkmark
\]

Since no two of $X$, $Y$, $Z$ commute, no two can be simultaneously definite.

% ============================================================
\section{2. Application: Two-Level System with System Matrix $H = \frac{\omega}{2}Z$}
% ============================================================

\subsection{2.1 The system matrix}
For a two-level system with energy splitting $\omega$:
\[
    H = \frac{\omega}{2}\,Z = \frac{\omega}{2}\begin{pmatrix}1&0\\0&-1\end{pmatrix}.
\]
The energy eigenstates and eigenvalues:
\[
    E_+ = +\frac{\omega}{2}\;(\mathbf{e}_0), \qquad
    E_- = -\frac{\omega}{2}\;(\mathbf{e}_1).
\]
The energy splitting is $\Delta E = \omega$.

\subsection{2.2 Time evolution}
Starting from a general state $\mathbf{v}(0)=\alpha\,\mathbf{e}_0+\beta\,\mathbf{e}_1$:
\[
    \mathbf{v}(t) = \alpha\,e^{-i\omega t/2}\,\mathbf{e}_0 + \beta\,e^{+i\omega t/2}\,\mathbf{e}_1.
\]
The magnitudes $|\alpha|$ and $|\beta|$ are constant; only the relative phase changes.

\subsection{2.3 Equations of motion for expectation values}
Using $\frac{d}{dt}\langle M \rangle = \frac{i}{\hbar}\langle [H,M] \rangle$ with $H=\frac{\omega}{2}Z$ (setting $\hbar=1$):
\begin{align}
    \frac{d}{dt}\langle X \rangle &= -i\cdot\frac{\omega}{2}\langle[X,Z]\rangle
    = -i\cdot\frac{\omega}{2}\cdot(-2i)\langle Y \rangle = -\omega\langle Y \rangle, \\
    \frac{d}{dt}\langle Y \rangle &= -i\cdot\frac{\omega}{2}\langle[Y,Z]\rangle
    = -i\cdot\frac{\omega}{2}\cdot(2i)\langle X \rangle = +\omega\langle X \rangle, \\
    \frac{d}{dt}\langle Z \rangle &= -i\cdot\frac{\omega}{2}\langle[Z,Z]\rangle = 0.
\end{align}
Notice that the factors of $i$ cancel, giving real equations for real expectation values.

\subsection{2.4 Precession}
The equations $\frac{d}{dt}\langle X \rangle=-\omega\langle Y \rangle$ and $\frac{d}{dt}\langle Y \rangle=+\omega\langle X \rangle$ describe \textbf{uniform circular motion} in the $XY$-plane at angular frequency $\omega$.  Combined with $\langle Z \rangle=\text{const}$:
\begin{align}
    \langle X \rangle(t) &= 2\operatorname{Re}\!\bigl(\alpha^*\beta\,e^{i\omega t}\bigr), \\
    \langle Y \rangle(t) &= 2\operatorname{Im}\!\bigl(\alpha^*\beta\,e^{i\omega t}\bigr), \\
    \langle Z \rangle(t) &= |\alpha|^2 - |\beta|^2.
\end{align}
The expectation-value vector $\bigl(\langle X \rangle,\langle Y \rangle,\langle Z \rangle\bigr)$ traces out a cone---this is \textbf{Larmor precession}, exactly matching the classical result for a magnetic dipole in a uniform field.

\subsection{2.5 General two-level system matrix}
The most general self-adjoint matrix on $\mathbb{C}^2$ (up to a constant multiple of $I$) is
\[
    H = \frac{\omega}{2}(n_x X + n_y Y + n_z Z),
\]
where $(n_x,n_y,n_z)$ is a unit vector.  The expectation values precess around this direction at frequency $\omega$, and the component along this direction is conserved.  The special case $(0,0,1)$ recovers the example above.

% ============================================================
\section{3. Summary}
% ============================================================

\begin{enumerate}[nosep]
    \item Two observables can be simultaneously known iff their matrices commute: $[L,M]=0$.
    \item Non-commuting observables (e.g.\ $X$ and $Z$) lead to \emph{uncertainty}: no state has definite values for both.
    \item The commutation relations are $[X,Y]=2iZ$ (and cyclic), verified by explicit matrix multiplication.
    \item A two-level system with $H=\frac{\omega}{2}Z$ has energy eigenstates $\mathbf{e}_0$ ($E=+\omega/2$) and $\mathbf{e}_1$ ($E=-\omega/2$).
    \item Expectation values evolve via $\frac{d}{dt}\langle L \rangle = \frac{i}{\hbar}\langle [H,L] \rangle$; for $H \propto Z$, $\langle X \rangle$ and $\langle Y \rangle$ precess while $\langle Z \rangle$ is conserved.
    \item The most general two-level system matrix is $H=\frac{\omega}{2}(n_x X + n_y Y + n_z Z)$.
\end{enumerate}

\end{document}
